% Slide — Portfólio final: 12 -> 18
\begin{frame}{Portfólio final: de 12 para 18 políticas}
  \centering
  \resizebox{\linewidth}{!}{%
  \tabstyle
  \begin{tabular}{@{}llll@{}}
    \toprule
    \textbf{Clássicas} & \textbf{SOTA / Zabraoui \novo} & \textbf{Busca \& RL} & \textbf{Híbridas} \\
    \midrule
    EOQ & PIL & GA & GA-DQN \\
    $(s,S)$ & CappedBaseStock & SA & GA-PPO \\
    Jornaleiro & BigDataNewsvendor & PSO & \\
     & Min-Max & DE & \\
     & Fixed Interval & DQN & \\
     & Vendor-Responsive & PPO & \\
     & & SARSA & \\
    \midrule
    \textbf{3} & \textbf{6 novas} & \textbf{7} & \textbf{2} \\
    \bottomrule
  \end{tabular}}
  \vspace{1em}

  \small \textbf{3 inalteradas} $\cdot$ \textbf{6 novas} $\cdot$
  \textbf{4 corrigidas} (GA, SA, PSO, DE por vetorização/aptidão; PPO por
  gradiente) $\cdot$ \textbf{5 reescritas} sem defeito conhecido (DQN,
  SARSA, GA-DQN, GA-PPO por herança das correções, simulador por
  vetorização).

  \note{
    O número final é 18: as 12 políticas de julho, menos zero (nenhuma foi
    removida), mais as 6 novas -- 3 do slot de estado da arte clássico
    (PIL, capped base-stock, big data newsvendor) e 3 do artigo-base
    Zabraoui (Min-Max, Fixed Interval, Vendor-Responsive). Cada célula
    deste portfólio tem agora uma justificativa documentada: ou é
    baseline preservado por comparabilidade, ou é a referência mais
    recente publicada para aquele slot conceitual, ou -- quando não existe
    referência publicada -- é a formulação canônica do método com a
    função objetivo corrigida.
  }
\end{frame}
